\chapter{Lie Groups}\label{lie-groups}
\chapterstatus{complete}

\section{Introduction}\label{lie-groups:section-introduction}

A Lie group is a group that is also a smooth manifold, with smooth group operations. Its local structure is captured by a finite-dimensional Lie algebra, and many of the spaces
of Hodge theory are homogeneous spaces of Lie groups: projective spaces, Grassmannians, flag manifolds, and the period domains that classify Hodge structures
(Chapter~\ref{vhs}). This chapter proves the correspondence between one-parameter subgroups and the Lie algebra, Cartan's closed subgroup theorem
(Theorem~\ref{lie-groups:theorem-closed-subgroup}) and the construction of homogeneous spaces (Theorem~\ref{lie-groups:theorem-homogeneous}). References:
\cite{knapp-lie}, \cite{brocker-tomdieck}, \cite{warner-foundations}.

\noindent\textbf{Prerequisites.} Chapter~\ref{smooth-manifolds} (Smooth Manifolds) and Chapter~\ref{representation-theory} (Representations of Groups and Lie Algebras).

\section{Lie groups and their Lie algebras}\label{lie-groups:section-lie-algebras}

\begin{definition}\label{lie-groups:definition-lie-group}
A \emph{Lie group} is a group $G$ with a smooth structure for which multiplication $G \times G \to G$ and inversion $G \to G$ are smooth. A \emph{homomorphism of Lie groups} is a smooth group
homomorphism. For $g \in G$, left and right translation are the diffeomorphisms $L_g(x) = gx$ and $R_g(x) = xg$.
\end{definition}

\begin{example}\label{lie-groups:example-lie-groups}
$(\RR^n, +)$; the torus $\RR^n/\ZZ^n$; $\GL_n(\RR)$ and $\GL_n(\CC)$, open subsets of matrix spaces, with polynomial multiplication and rational inversion; the \emph{classical groups}
$\SL_n$, $O(n) = \{A : A^TA = I\}$, $\mathrm{SO}(n)$, $U(n) = \{A : A^*A = I\}$ and $\mathrm{SU}(n)$, which are submanifolds by Example~\ref{smooth-manifolds:example-regular-value} and subgroups, hence Lie
groups. Products of Lie groups are Lie groups, and the identity component $G^0$ is an open normal subgroup.
\end{example}

\begin{definition}\label{lie-groups:definition-lie-algebra-of-g}
A vector field $X$ on $G$ is \emph{left-invariant} if $d(L_g)X = X \circ L_g$ for all $g$. The \emph{Lie algebra} of $G$ is $\mathfrak{g} = T_eG$, with the bracket transported from the left-invariant vector
fields.
\end{definition}

\begin{proposition}\label{lie-groups:proposition-left-invariant}
The map $X \mapsto X_e$ is a bijection from left-invariant vector fields to $T_eG$, its inverse sending $v$ to the field $g \mapsto d(L_g)_ev$; left-invariant fields are smooth and closed under the Lie
bracket. Hence $\mathfrak{g}$ is a Lie algebra of dimension $\dim G$ (Definition~\ref{representation-theory:definition-lie-algebra}), and a homomorphism $\varphi \colon G \to H$ induces a homomorphism of Lie
algebras $d_e\varphi \colon \mathfrak{g} \to \mathfrak{h}$.
\end{proposition}

\begin{proof}
Left-invariance determines the field from its value at $e$, and $g \mapsto d(L_g)_ev$ is smooth because it is the image of $(g, v)$ under the smooth map $G \times \mathfrak{g} \to TG$ induced by
multiplication. If $X$, $Y$ are left-invariant, so is $[X, Y]$, since $d(L_g)$ preserves brackets: for a diffeomorphism $F$, $[F_*X, F_*Y] = F_*[X, Y]$, as both sides act on functions by
$f \mapsto (X(Y(f \circ F)) - Y(X(f \circ F))) \circ F^{-1}$. For the last statement, $\varphi \circ L_g = L_{\varphi(g)} \circ \varphi$ shows that $d_e\varphi$ intertwines the corresponding left-invariant fields,
which are therefore $\varphi$-related; and brackets of $\varphi$-related fields are $\varphi$-related.
\end{proof}

\begin{example}\label{lie-groups:example-matrix-lie-algebras}
For $G \subseteq \GL_n$ a matrix group, $\mathfrak{g} \subseteq M_n$ and the bracket is the commutator $XY - YX$: indeed for $G = \GL_n$ the left-invariant field with value $X$ at $I$ is $A \mapsto AX$, and a
computation in the coordinates of $M_n$ gives $[X, Y] = XY - YX$. Thus $\mathfrak{gl}_n = M_n$, $\mathfrak{sl}_n = \{X : \operatorname{tr}X = 0\}$, $\mathfrak{o}(n) = \{X : X^T = -X\}$,
$\mathfrak{u}(n) = \{X : X^* = -X\}$ and $\mathfrak{su}(n) = \mathfrak{u}(n) \cap \mathfrak{sl}_n(\CC)$, by differentiating the defining equations at $I$.
\end{example}

\section{The exponential map}\label{lie-groups:section-exponential}

\begin{theorem}\label{lie-groups:theorem-exponential}
For $X \in \mathfrak{g}$ the left-invariant field $X$ is complete, and its flow through $e$ is a homomorphism $\gamma_X \colon \RR \to G$; conversely every homomorphism $\RR \to G$ arises this way. Setting
$\exp(X) = \gamma_X(1)$ defines a smooth map $\exp \colon \mathfrak{g} \to G$ with $d_0\exp = \id$, so $\exp$ is a diffeomorphism from a neighbourhood of $0$ onto a neighbourhood of $e$. It satisfies
$\exp((s + t)X) = \exp(sX)\exp(tX)$, $\exp(-X) = \exp(X)^{-1}$, and $\varphi(\exp X) = \exp(d_e\varphi(X))$ for every homomorphism $\varphi$. For matrix groups, $\exp(X) = \sum_kX^k/k!$.
\end{theorem}

\begin{proof}
Let $\gamma$ be the integral curve of $X$ with $\gamma(0) = e$, defined on $(-\varepsilon, \varepsilon)$. For fixed $s$, both $t \mapsto \gamma(s)\gamma(t)$ and $t \mapsto \gamma(s + t)$ are integral curves through
$\gamma(s)$ (left-invariance), so they agree where defined; this extends $\gamma$ to all of $\RR$ and gives the homomorphism property. Conversely, a homomorphism $\gamma \colon \RR \to G$ is the integral curve of
the left-invariant field with value $\gamma'(0)$, by the same computation. Smoothness of $(t, X) \mapsto \gamma_X(t)$ follows from smooth dependence of solutions of differential equations on parameters
(Theorem~\ref{real-analysis:theorem-smooth-dependence}), applied to the field on $G \times \mathfrak{g}$ whose value at $(g, X)$ is $(d(L_g)X, 0)$. Since $\gamma_{sX}(t) = \gamma_X(st)$ we get
$\exp(sX) = \gamma_X(s)$, hence $d_0\exp(X) = \frac{d}{dt}\big|_0\exp(tX) = X$ and the group laws; the inverse function theorem gives the local diffeomorphism. Naturality holds because
$t \mapsto \varphi(\exp tX)$ is a homomorphism $\RR \to H$ with derivative $d_e\varphi(X)$ at $0$. For matrix groups, $t \mapsto \sum_kt^kX^k/k!$ converges (Proposition~\ref{real-analysis:proposition-series}),
is a homomorphism by the Cauchy product, and has derivative $X$ at $t = 0$.
\end{proof}

\begin{definition}\label{lie-groups:definition-adjoint}
The \emph{adjoint representation} is $\Ad \colon G \to \GL(\mathfrak{g})$, $\Ad(g) = d_e(c_g)$ where $c_g(x) = gxg^{-1}$; its differential is $\ad = d_e\Ad \colon \mathfrak{g} \to \mathfrak{gl}(\mathfrak{g})$.
\end{definition}

\begin{proposition}\label{lie-groups:proposition-adjoint}
$\Ad$ is a homomorphism of Lie groups, $\ad(X)(Y) = [X, Y]$, and $\Ad(\exp X) = e^{\ad X}$. Moreover $\exp(X)\exp(Y) = \exp(X + Y + \frac12[X, Y] + \cdots)$ for $X$, $Y$ small, where the omitted terms are
iterated brackets (the Baker--Campbell--Hausdorff formula \cite[Chapter 1]{knapp-lie}).
\end{proposition}

\begin{proof}
$\Ad$ is smooth because $(g, x) \mapsto gxg^{-1}$ is, and it is a homomorphism since $c_{gh} = c_g \circ c_h$. For the bracket, note that $c_g$ is a Lie group automorphism, so $\Ad(g)$ preserves brackets; and
$\frac{d}{dt}\big|_0\Ad(\exp tX)Y$ is by definition the Lie derivative of $Y$ along $X$, which is $[X, Y]$ (Proposition~\ref{smooth-manifolds:proposition-bracket}), since the flow of the left-invariant field
$X$ is right translation by $\exp(tX)$ and $\Ad(\exp tX)$ is the induced map on left-invariant fields. Finally $t \mapsto \Ad(\exp tX)$ is a homomorphism $\RR \to \GL(\mathfrak{g})$ with derivative $\ad X$, hence
equals $e^{t\,\ad X}$ by Theorem~\ref{lie-groups:theorem-exponential}.
\end{proof}

\section{Subgroups and homogeneous spaces}\label{lie-groups:section-homogeneous}

\begin{theorem}[Cartan]\label{lie-groups:theorem-closed-subgroup}
A closed subgroup $H$ of a Lie group $G$ is an embedded submanifold, hence a Lie group, with Lie algebra
$\mathfrak{h} = \{X \in \mathfrak{g} : \exp(tX) \in H \text{ for all } t \in \RR\}$.
\end{theorem}

\begin{proof}
Let $\mathfrak{h}$ be as stated. It is closed under scalars by definition. It is closed under addition and brackets: by the product formulas
$\exp(t(X + Y)) = \lim_k(\exp(tX/k)\exp(tY/k))^k$ and $\exp(t^2[X, Y]) = \lim_k(\exp(tX/k)\exp(tY/k)\exp(-tX/k)\exp(-tY/k))^{k^2}$, which follow from the Baker--Campbell--Hausdorff formula or from a direct
Taylor expansion, and $H$ is closed. So $\mathfrak{h}$ is a Lie subalgebra.

Choose a complement $\mathfrak{m}$ with $\mathfrak{g} = \mathfrak{h} \oplus \mathfrak{m}$. We claim that $\exp(\mathfrak{m})$ meets $H$ only at $e$ near $0$: otherwise there are $Y_k \in \mathfrak{m} \setminus \{0\}$
with $Y_k \to 0$ and $\exp Y_k \in H$; passing to a subsequence, $Y_k/\|Y_k\| \to Y \in \mathfrak{m}$ with $\|Y\| = 1$. For $t \in \RR$ let $m_k = \lfloor t/\|Y_k\|\rfloor$; then
$m_k Y_k \to tY$ and $\exp(m_kY_k) = (\exp Y_k)^{m_k} \in H$, so $\exp(tY) \in H$ by closedness, giving $Y \in \mathfrak{h} \cap \mathfrak{m} = 0$, a contradiction.

The map $\mathfrak{h} \times \mathfrak{m} \to G$, $(X, Y) \mapsto \exp X\exp Y$, has invertible differential at $(0,0)$, hence is a diffeomorphism near $0$; by the claim, in a small enough neighbourhood it
carries $\mathfrak{h} \times \{0\}$ onto $H$. Translating by elements of $H$ gives charts of $G$ in which $H$ is a slice, so $H$ is an embedded submanifold; the group operations are smooth as restrictions,
and $T_eH = \mathfrak{h}$.
\end{proof}

\begin{theorem}\label{lie-groups:theorem-homogeneous}
Let $H \subseteq G$ be a closed subgroup. The quotient $G/H$ has a unique smooth structure such that the projection $\pi \colon G \to G/H$ is a submersion; it is a manifold of dimension $\dim G - \dim H$ on
which $G$ acts smoothly and transitively. If a Lie group $G$ acts smoothly and transitively on a manifold $M$ and $H$ is the stabiliser of $p \in M$, then $G/H \to M$, $gH \mapsto gp$, is a
$G$-equivariant diffeomorphism.
\end{theorem}

\begin{proof}
With $\mathfrak{m}$ as in the previous proof, the map $\mathfrak{m} \to G/H$, $Y \mapsto \pi(\exp Y)$, is injective near $0$ and its image is open; these maps, translated by $G$, are the charts. The transition
maps are smooth because multiplication in $G$ is, and $\pi$ becomes the projection $(X, Y) \mapsto Y$ in the corresponding charts, so it is a submersion; uniqueness follows because a submersion determines the
smooth structure on its image via local sections. The action $G \times G/H \to G/H$ is smooth because it is so after composing with the submersion $\id \times \pi$. For the last statement, consider the orbit map $f \colon G \to M$, $g \mapsto gp$, which is smooth and surjective. If $df$ were nowhere surjective, every point of $M$ would be a
critical value, and $M$ would have measure zero in itself by Sard's theorem (Theorem~\ref{smooth-manifolds:theorem-sard}), which is false. So $df$ is surjective at some $g_0$, hence everywhere, because
$f \circ L_g = g \cdot f$ and the action of $g$ is a diffeomorphism of $M$. Thus $f$ is a submersion, and so is the induced smooth bijection $G/H \to M$ (as $\pi$ is a surjective submersion); a bijective
submersion is a diffeomorphism, by the inverse function theorem and equality of dimensions.
\end{proof}

\begin{example}\label{lie-groups:example-homogeneous}
\begin{enumerate}
\item $S^{n} = O(n+1)/O(n) = \mathrm{SO}(n+1)/\mathrm{SO}(n)$ and $S^{2n+1} = U(n+1)/U(n)$.
\item $\RR\PP^n = O(n+1)/(O(n) \times O(1))$ and $\CC\PP^n = U(n+1)/(U(n) \times U(1))$; more generally the Grassmannian $\mathrm{Gr}(k, n)$ of $k$-planes in $\RR^n$ is $O(n)/(O(k) \times O(n-k))$, a compact
manifold of dimension $k(n - k)$, and similarly over $\CC$.
\item Flag manifolds $U(n)/(U(n_1) \times \cdots \times U(n_r))$; these and their open subsets are the period domains parametrising Hodge structures (Chapter~\ref{vhs}).
\item $\GL_n(\RR)/O(n)$ is the space of positive definite symmetric matrices, and the upper half plane is $\SL_2(\RR)/\mathrm{SO}(2)$, the period domain for Hodge structures of weight one and rank two
(Chapter~\ref{hodge-structures}).
\end{enumerate}
\end{example}

\section{Compact Lie groups}\label{lie-groups:section-compact}

\begin{theorem}[Haar measure]\label{lie-groups:theorem-haar}
On a compact Lie group $G$ there is a unique regular Borel probability measure $\mu$ that is invariant under all left translations; it is also right invariant and invariant under inversion. Consequently
$\int_Gf(gx)\,d\mu(x) = \int_Gf(x)\,d\mu(x)$ for all $g$ and all continuous $f$.
\end{theorem}

This is \cite[Chapter I]{brocker-tomdieck}; the measure can be constructed by integrating a left-invariant volume form (Chapter~\ref{differential-forms}), which exists because
$\Lambda^{\dim G}\mathfrak{g}^*$ is one-dimensional, and normalising. Uniqueness follows by convolving two invariant measures.

\begin{corollary}\label{lie-groups:corollary-averaging}
Every finite-dimensional representation of a compact Lie group $G$ on a complex vector space admits a $G$-invariant Hermitian inner product, and is a direct sum of irreducible subrepresentations.
\end{corollary}

\begin{proof}
Average any Hermitian inner product: $\langle v, w \rangle_G = \int_G\langle gv, gw \rangle\,d\mu(g)$, which is Hermitian, positive definite and invariant. The orthogonal complement of a subrepresentation is
then a subrepresentation, and induction on the dimension gives the decomposition, exactly as in Maschke's theorem (Theorem~\ref{representation-theory:theorem-maschke}).
\end{proof}

\begin{remark}\label{lie-groups:remark-structure}
Deeper structure theory, which we quote when needed: a connected Lie group has a unique simply connected covering group, and Lie algebra homomorphisms integrate to group homomorphisms from the simply
connected group (Lie's second theorem); Lie subalgebras correspond to connected Lie subgroups; every compact connected Lie group is a quotient of a product of a torus and simply connected compact simple
groups, which are classified by their Dynkin diagrams; and every connected compact Lie group is the maximal compact subgroup of a unique complex reductive group, its \emph{complexification} (for instance
$U(n) \subseteq \GL_n(\CC)$). See \cite{knapp-lie} and \cite{brocker-tomdieck}. The unitary trick, which deduces complete reducibility of representations of complex reductive groups from
Corollary~\ref{lie-groups:corollary-averaging}, is used in Chapter~\ref{mumford-tate}.
\end{remark}

\section{Exercises}\label{lie-groups:section-exercises}

\begin{exercise}\label{lie-groups:exercise-exp-not-surjective}
Show that $\exp \colon \mathfrak{sl}_2(\RR) \to \SL_2(\RR)$ is not surjective, although $\exp$ is surjective for compact connected Lie groups and for $\GL_n(\CC)$.
\end{exercise}

\begin{solution}
An element of the image with distinct real eigenvalues has positive eigenvalues, and $\operatorname{diag}(-2, -1/2) \in \SL_2(\RR)$ has negative ones; it is not $\exp$ of a real matrix, since the
eigenvalues of $\exp X$ are exponentials of those of $X$ and complex eigenvalues of a real matrix come in conjugate pairs.
\end{solution}

\begin{exercise}\label{lie-groups:exercise-su2}
Show that $\mathrm{SU}(2)$ is diffeomorphic to $S^3$, that its Lie algebra is $\mathfrak{su}(2) \cong \RR^3$ with the cross product, and that the adjoint representation gives a double covering
$\mathrm{SU}(2) \to \mathrm{SO}(3)$.
\end{exercise}

\begin{exercise}\label{lie-groups:exercise-closed-subgroup-needed}
Show that the image of the homomorphism $\RR \to T^2$, $t \mapsto (t, \alpha t) \bmod \ZZ^2$, with $\alpha$ irrational, is a dense subgroup that is not closed and not an embedded submanifold; so
``closed'' cannot be dropped in Theorem~\ref{lie-groups:theorem-closed-subgroup}.
\end{exercise}

\begin{exercise}\label{lie-groups:exercise-grassmannian-dimension}
Verify that $\mathrm{Gr}(k, n) = O(n)/(O(k) \times O(n-k))$ has dimension $k(n-k)$ and is connected, and identify $\mathrm{Gr}(1, n)$ with $\RR\PP^{n-1}$.
\end{exercise}
